\usepackage{amsthm,amssymb}
\usepackage{mathtools}
\usepackage[a4paper, margin=1in]{geometry}
\usepackage{float}
\usepackage{indentfirst}
\usepackage{anyfontsize}
\usepackage{booktabs,multirow,multicol}
\usepackage[shortlabels,inline]{enumitem}
\usepackage{appendix}
\usepackage[UTF8]{ctex}
\usepackage{fontspec}
\usepackage{booktabs}

\geometry{a4paper,left=1.5cm,right=1.5cm,top=1.5cm,bottom=1.5cm}
\setmainfont{Times New Roman} % 设置正文主字体
\setsansfont{Arial} % 设置无衬线字体
\setmonofont{Courier New} % 设置等宽字体
\setCJKmainfont{SimSun}[AutoFakeBold=2.0] % 设置中文主字体（宋体）
\setCJKsansfont{SimHei} % 设置中文无衬线字体（黑体）
\setCJKmonofont{FangSong} % 设置中文等宽字体（仿宋）
\pagestyle{plain}
\setcounter{tocdepth}{1}

\usepackage[dvipsnames,svgnames]{xcolor}
\usepackage{graphicx}
\graphicspath{
    {./figure/}{./figures/}{./image/}{./images/}{./graphic/}{./graphics/}{./picture/}{./pictures/}
}
\usepackage{subcaption}

\usepackage[ruled,linesnumbered,noline]{algorithm2e}
\usepackage{listings}
\lstdefinestyle{simpleStyle}{
    basicstyle=\ttfamily\small,
    breaklines=true,
    keywordstyle=\color{blue},
    identifierstyle=\color{black},
    stringstyle=\color{violet},
    commentstyle=\color[RGB]{34,139,34},
    showstringspaces=false,
    numbers=left,
    numbersep=2em,
    numberstyle=\footnotesize,
    frame=single,
    framesep=1em,
}
\lstset{style=simpleStyle}

\usepackage{hyperref}
\hypersetup{
    colorlinks=true,linkcolor=,urlcolor=cyan
}

\usepackage{thmtools}

%% define environments

\declaretheorem[style=plain, name=Lexical Discrimination, numbered=no]{lex}

\declaretheorem[style=definition, name=Word Origin, numbered=no]{origin}

\declaretheorem[style=definition, name=Language Focus, numbered=no]{gram}

\usepackage[most]{tcolorbox}

\newcommand{\newtcbenvironment}[2]{
    \tcolorboxenvironment{#1}{#2, enhanced, breakable, sharp corners, boxrule=0pt, colframe=white}
    \tcolorboxenvironment{#1*}{#2, enhanced, breakable, rounded corners, boxrule=0pt, colframe=white}
}

%% define styles

\newtcbenvironment{lex}{colframe=RoyalPurple, colback=RoyalPurple!8}
\newtcbenvironment{origin}{colframe=ForestGreen, colback=ForestGreen!5}
\newtcbenvironment{gram}{colframe=RawSienna, colback=RawSienna!5}

%% cbox
\newtcolorbox{cbox}[1][]{%
    enhanced,
    breakable,
    sharp corners,
    leftrule=2pt, rightrule=0pt, toprule=0pt, bottomrule=0pt,
    colframe=SkyBlue,
    colback=SkyBlue!8,
    #1
}

\usepackage{pifont}

% 词条标题: \vw{单词}{音标}{词性速览}
\newcommand{\vw}[3]{%
  \par\bigskip\noindent
  {\color{SkyBlue!70!NavyBlue}\bfseries\Large #1}\
  {\small [{#3}]}\
  {\small\color{gray} {#2}}%
  \par\smallskip
}

% 短语词条标题: \vp{短语}
\newcommand{\vp}[1]{%
  \par\bigskip\noindent
  {\color{SkyBlue!70!NavyBlue}\bfseries\Large #1}%
  \par\smallskip
}

% 义项: \sense{词性}{中文释义}
\definecolor{mygreen}{RGB}{53, 94, 59}
\newcommand{\sense}[2]{\bfseries\par\medskip\noindent{\color{mygreen}\bfseries[{#1}]}\ {\color{mygreen}\bfseries {#2}}\par\smallskip}

% 搭配/用法行: \col{搭配或句型}
\definecolor{myBurgundy}{RGB}{128,0,32}
\newcommand{\col}[1]{\par\noindent\hspace{1em}{\color{myBurgundy}$\blacktriangleright$\ \bfseries {#1}}\par\nopagebreak\smallskip}

% 例句(广义,含短语级示例): \ex{英文}
\newcommand{\ex}[1]{\par\nopagebreak\noindent\hspace{2.2em}{{#1}}\par\smallskip}
